\subsection{Compresión de Memoria sin Pérdida de Precisión}
La simulación cuántica eficiente requiere reducir el consumo de memoria sin
comprometer la fidelidad del estado cuántico. Para ello, se han desarrollado
técnicas de compresión sin pérdida de precisión que buscan representar estados
y operadores cuánticos de manera más compacta sin alterar sus valores.

Uno de los primeros enfoques en esta dirección fue el uso de diagramas de decisión
cuánticos (\textit{Quantum Information Decision Diagrams}, QuIDD),
introducidos por \textcite{Viamontes2003}. Este método permite representar
vectores de estado y operadores unitarios aprovechando redundancias
estructurales en sus coeficientes, lo que permite una representación más eficiente
en memoria. Su ventaja radica en que, para circuitos con alta redundancia, el
crecimiento en memoria puede reducirse de exponencial a polinomial. Sin embargo,
su eficacia se limita a circuitos con estructura repetitiva, fallando en casos
de entrelazamiento complejo.

En un desarrollo posterior, \textcite{Miller2006} propusieron los \textit{Quantum
Multiple-valued Decision Diagrams} (QMDD), los cuales extienden la
representación de QuIDD mediante el uso de aristas ponderadas con matrices
unitarias. Esto permitió manejar de manera más eficiente circuitos reversibles
y operadores cuánticos de mayor tamaño. Aunque estos diagramas optimizan la
representación de puertas lógicas y matrices densas, su eficiencia sigue dependiendo
de la estructura interna del circuito.

Más recientemente, \textcite{Jaques2022} exploraron la esparsidad del estado
cuántico para almacenar solo las amplitudes no nulas, reduciendo significativamente
el almacenamiento requerido en algoritmos con exploración limitada del espacio
de Hilbert. Su enfoque demostró ser altamente eficiente en problemas de
factorización y aritmética cuántica, pero pierde efectividad en circuitos con alta
interferencia y entrelazamiento.

En 2023, \textcite{Vinkhuijzen2023} introdujeron los \textit{Local Invertible
Map Decision Diagrams} (LIMDD), una evolución de los diagramas de decisión que
incorpora transformaciones locales para mejorar la representación compacta del
estado cuántico. Esta técnica permite manejar estados estabilizadores y ciertas
estructuras altamente entrelazadas que eran difíciles de representar con diagramas
de decisión tradicionales. No obstante, su aplicabilidad sigue limitada a casos
donde las transformaciones locales puedan ser explotadas para reducir la complejidad
estructural.

Estos avances en compresión sin pérdida han permitido extender la capacidad de
simulación de circuitos cuánticos en hardware clásico. Sin embargo, conviene distinguir
tres familias de soluciones. Unas cambian la representación matemática del estado, como
los diagramas de decisión; otras conservan exactitud, pero dependen de estructura
especial, por ejemplo esparsidad o simetrías explotables; y otras todavía dejan abierto el
problema de reducir memoria dentro del formalismo Schr\"odinger con almacenamiento
explícito del vector de amplitudes.

Esa distinción es precisamente la brecha que motiva esta tesis. El espacio menos cubierto
por la literatura revisada no es el de representar el estado de otra manera, sino el de
comprimir sin pérdida un vector de estado explícito, manteniendo la semántica del
simulador tipo Schr\"odinger y gestionando la sobrecarga asociada a acceso, caché y
actualización de bloques.
